% ===================================================================
%  TABEL Nr. 6 — Maja seest, mööbel, toitumine, valgustus.
%  Traduction 2026 d'apres texto/io, controlee sur texto/fr et
%  texto/fi. Consigne : voir texto/et/10-tabel-01.tex.
%
%  LE TABLEAU 6 ROMANCHE AVAIT MONTRE QU'UNE TRADUCTION DIT AUSSI CE
%  QUE LE LECTEUR N'A PAS CHEZ LUI ; L'ESTONIEN LE DIT PAR LA MEME
%  PAGE ET PAR UNE AUTRE PIECE.
%
%  Le foyer de cette maison est estonien de bout en bout : « ahi » le
%  poele, « tuli » le feu, « kolle » l'atre, « pott », « pada »,
%  « kulp » la louche, « sõel » le tamis, « leib » le pain, « või »
%  le beurre, « piim » le lait, « sool » le sel. Pas un emprunt : ce
%  sont les mots d'avant les seigneurs.
%
%  La salle de bains de la derniere section ne l'est pas d'un mot :
%  « vannituba » la piece-a-baignoire, « dušš », « kraan »,
%  « kahhelplaat ». C'est une installation de ville du XIXe siecle,
%  arrivee avec l'allemand et son plombier.
%
%  ET LE MOT QUI MANQUE EST CELUI QUE L'EUROPE ENTIERE A PRIS. La
%  seule piece qu'un Estonien lave son corps depuis mille ans porte
%  un nom balto-finnois, « saun » — le meme que le finnois a donne au
%  monde entier sous la forme « sauna ». C'est le seul mot de cette
%  famille de langues que le francais, l'anglais, l'allemand et le
%  russe ecrivent sans le traduire. ET LA PLANCHE NE LE DEMANDE PAS :
%  elle montre une baignoire, une douche et un bassin a pieds.
%
%  LE ROMANCHE AVAIT DU EMPLOYER SON MOT LE PLUS RICHE, « pigna »,
%  POUR LA SEULE LIGNE OU LA PLANCHE PARLE DE POELE. L'estonien ne
%  peut meme pas employer le sien : la planche ne montre pas la piece
%  ou il servirait. Deux langues, deux facons de tomber a cote, et la
%  meme cause — LA PLANCHE EST LA MAISON D'UN AUTRE.
%
%  RESTITUTION D'UNE PARENTHESE, BLOC t06-c3-08-2. Le fac-simile ido
%  ecrit « « La Hotelo » n° 13) » sans la parenthese ouvrante. La
%  coquille est du compositeur de l'ido : la parenthese est retablie,
%  comme dans la colonne romanche. UNE TRADUCTION N'EST PAS UNE
%  TRANSCRIPTION.
% ===================================================================

%%K t06-apar-1 apar
\VUsaut{6.0mm}
\VUtitre{15.0pt}{60}{TABEL Nr. 6}
\VUsaut{1.4mm}
\VUtitre{11.4pt}{0}{Maja seest, mööbel, toitumine,}
\VUsaut{0.4mm}
\VUtitre{11.4pt}{0}{valgustus.}
\VUsaut{2.2mm}

%%K t06-apar-2 apar
Maja on valmis. Jaani onu on sisse seadnud oma uues kodus, mille korralduse me juba
tunneme. Nüüd vaadakem, kuidas ta on oma maja sisustanud. Me uurime kõigepealt:

%%K t06-tit-1 sub
\VUpk{10.2pt}{40}{Magamistuba.}
\VUsaut{3.0mm}

%%K t06-c1-01-1 p
1. --- Me tuleme paraja hetkel, sest toatüdruk on ametis \VUgras{toa} koristamisega.

%%K t06-c1-02-1 p
2. --- \VUgras{Tualettlaual} \textsuperscript{(1)} on siin ja seal mitmesuguseid
esemeid, millega pestakse, kammitakse, lühidalt tehakse tualetti. Need on
\VUgras{kauss} \textsuperscript{(2)} ja \VUgras{veekann} \textsuperscript{(3)},
\VUgras{juuksehari} \textsuperscript{(4)}, \VUgras{kamm} \textsuperscript{(5)},
\VUgras{seep} \textsuperscript{(6)} ja \VUgras{käsn} \textsuperscript{(7)}, lõpuks
\VUgras{pudeleke} \textsuperscript{(8)} vedela hambavee jaoks.

%%K t06-c1-03-1 p
3. --- Põrandal, tualettlaua ees, on \VUgras{ämber} \textsuperscript{(9)} musta vee
kogumiseks.

%%K t06-c1-04-1 p
4. --- \VUgras{Eestoas} \textsuperscript{(10)} näeme mõnda \VUgras{riidenagi}
\textsuperscript{(13)} ja kaht \VUgras{vihmavarju} \textsuperscript{(11)}
\VUgras{vihmavarjunõus} \textsuperscript{(12)}.

%%K t06-c1-05-1 p
5. --- Ukse teisel pool on \VUgras{peegelkapp} \textsuperscript{(14)}, mis sisaldab
\VUgras{pesu} \textsuperscript{(15)}.

%%K t06-c1-06-1 p
6. --- Kasutagem ära, et \VUgras{toatüdruk} \textsuperscript{(6)} teeb \VUgras{voodit}
\textsuperscript{(17)}, et uurida selle eri osi. \VUgras{Voodiraamis}
\textsuperscript{(20)} on hea \VUgras{vedrumadrats} \textsuperscript{(21)}, millele
Justiina kohe paneb villase \VUgras{madratsi} \textsuperscript{(22)}. Siis võtab ta
toolidelt kaks \VUgras{lina} \textsuperscript{(23)} ja \VUgras{teki}
\textsuperscript{(24)}. Seejärel paneb ta voodi peatsisse \VUgras{põikpadja}
\textsuperscript{(25)} ja \VUgras{peapadja} \textsuperscript{(26)}, mis on koos
\VUgras{sulgtekiga} \textsuperscript{(27)} \VUgras{voodikattel} \textsuperscript{(32)}.
Ta kasutab sulgtolmulappi, et tolmust puhastada \VUgras{kardinad}
\textsuperscript{(19)} ja \VUgras{voodi baldahhiin} \textsuperscript{(18)}, ja vaat et
osa tööst ongi tehtud.

%%K t06-c1-07-1 p
7. --- Siis peab ta pisut korrastama \VUgras{öökapi} \textsuperscript{(28)}, üles
keerama \VUgras{äratuskella} \textsuperscript{(29)}, ette valmistama \VUgras{öölambi}
\textsuperscript{(30)}, ära võtma \VUgras{küünla} \textsuperscript{(31)}, et küünlajalga
puhastada.

%%K t06-tit-2 sub
\VUsaut{5.0mm}
\VUpk{10.2pt}{40}{Töökorv ja kaminatarbed.}
\VUsaut{3.0mm}

%%K t06-c1-08-1 p
8. --- \VUgras{Töökorv} \textsuperscript{(34)} \VUgras{tabureti} \textsuperscript{(33)}
kõrval vajab samuti väga korrastamist. Ta peab kokku voltima \VUgras{tikandid}
\textsuperscript{(35)}, sulgema \VUgras{töölauda} \textsuperscript{(38)}
\VUgras{sõrmkübara}, \VUgras{niidi} \textsuperscript{(36)}, \VUgras{nõelad}
\textsuperscript{(37)} ja \VUgras{käärid} \textsuperscript{(39)} ning lukustama
\VUgras{võtmega} \textsuperscript{(40)} laua kaane.

%%K t06-c1-09-1 p
9. --- Keskel oleval \VUgras{ühe jalaga laual} \textsuperscript{(41)} vahetab Justiina
\VUgras{karahvini} \textsuperscript{(42)} vee. Ta paneb \VUgras{suhkrut}
\VUgras{suhkrutoosi} \textsuperscript{(43)}, puhastab \VUgras{suhkrutangid}
\textsuperscript{(44)} ja võtab ära \VUgras{riideharja} \textsuperscript{(46)}.

%%K t06-c1-10-1 p
10. --- Toa parem pool nõuab temalt vähem tööd, sest \VUgras{lamamistool}
\textsuperscript{(47)} ja tema \VUgras{padi} \textsuperscript{(48)} on oma õigel kohal,
ja kuna ilm ei ole külm, ei ole miski paigast nihkunud \VUgras{kamina}
\textsuperscript{(49)} tarvete seas. \VUgras{Labidas} \textsuperscript{(50)},
\VUgras{tangid} \textsuperscript{(51)}, \VUgras{tulekoerad} \textsuperscript{(55)},
\VUgras{tuhatõke} \textsuperscript{(56)} on puhtad; \VUgras{tulesirm}
\textsuperscript{(54)} ei ole paigast nihutatud ja \VUgras{puukast}
\textsuperscript{(52)} on hästi varustatud \VUgras{halgudega} \textsuperscript{(53)}.

%%K t06-noto-1 noto
\VUnotes{143.2mm}{%
(*) Babo = väikelaps, ingl. baby, pr. bébé, it. bambino.
\par
%%K t06-noto-2 noto
(*) Lambrequino = pr. lambrequin, hisp. lambrequines, port. lambrequins, isegi sks.
ingl. lambrequins, seega sks. ingl. pr. port. hisp.%
}

%%K t06-c1-11-1 p
11. --- Siiski peab ta tolmust puhastama \VUgras{pendelkella} \textsuperscript{(57)},
\VUgras{kandelaabrid} \textsuperscript{(58)} ja \VUgras{taskuasjade alused}
\textsuperscript{(59)}, lõpuks pühkima \VUgras{peeglit} \textsuperscript{(60)}.

%%K t06-c1-12-1 p
12. --- Mis puutub väikelapse (\VUgras{babo} (*)) \VUgras{hälli}
\textsuperscript{(61)}, siis see on juba valmis.

%%K t06-tit-3 sub
\VUsaut{5.0mm}
\VUpk{10.2pt}{40}{Salong.}
\VUsaut{3.0mm}

%%K t06-c2-01-1 p
1. --- Nüüd on siin salong. See on väga ilus \VUgras{tuba}. Kamin, mis asub paremas
nurgas, on ehitud õrna \VUgras{kujukesega} \textsuperscript{(1)}, kahe ilusa
\VUgras{vaasiga} \textsuperscript{(3)} ja drapeeritud sametist \VUgras{lambrekiiniga}
\textsuperscript{(2)} (*).

%%K t06-c2-02-1 p
2. --- Et takistada kolde sädemetel vaipa põletamast, on kolde ette pandud vasest
\VUgras{sädemekaitse} \textsuperscript{(4)}.

%%K t06-c2-03-1 p
3. --- Lähedal on avatud elegantne daami \VUgras{kirjutuslaud} \textsuperscript{(5)}.
Selle taga kirjutab \VUgras{perenaine} \textsuperscript{(23)} oma kirju.

%%K t06-c2-04-1 p
4. --- Aken on ehitud \VUgras{aknakardinatega} \textsuperscript{(6)}, mille
\VUgras{sidemed} \textsuperscript{(8)} on kinnitatud \VUgras{nagide}
\textsuperscript{(7)} külge, ja riidest suurte kardinatega, mis on \VUgras{ülalt
drapeeritud} \textsuperscript{(9)}.

%%K t06-c2-05-1 p
5. --- Akna süvendis on \VUgras{potihoidjas} \textsuperscript{(11)} roheline
\VUgras{taim} \textsuperscript{(10)}, uhke palm, mille lehed ulatuvad peaaegu
\VUgras{petrooleumilambini} \textsuperscript{(12)}.

%%K t06-c2-06-1 p
6. --- Lambi \VUgras{lambivarju} \textsuperscript{(13)} on seadnud Maria, võluv
\VUgras{neiu} \textsuperscript{(14)}, kes on klaveri \textsuperscript{(15)} juures. Väga
sirgelt \VUgras{pingil} \textsuperscript{(16)} istudes õpib ta muusikapala. Ta on väga
osav ja hoolikas, mida tõendab tema \VUgras{noodikapp} \textsuperscript{(17)}, mis on
täiuslikus korras.

%%K t06-tit-4 sub
\VUsaut{5.0mm}
\VUpk{10.2pt}{40}{Üleannetu poiss.}
\VUsaut{3.0mm}

%%K t06-c2-07-1 p
7. --- Tema vend Paul otsib \VUgras{raamatut} \textsuperscript{(19)}
\VUgras{raamatukapist} \textsuperscript{(18)}. Ta on \VUgras{noormees}
\textsuperscript{(20)}, rahutu ja talumatu. Kapi külge klammerdudes, et ulatuda
raamatuni, mis on liiga kõrgel, riskib ta maha ajada \VUgras{büsti}
\textsuperscript{(21)}, mis on ülal.

%%K t06-c2-08-1 p
8. --- Ta kasutab selle rumaluse tegemiseks ära, et tema ema on ametis vestlemisega
sõbratariga, \VUgras{daamiga} \textsuperscript{(22)}, kes tuli mõneks päevaks tema koju.
Ema istub \VUgras{diivanil} \textsuperscript{(24)} \VUgras{tugitooli}
\textsuperscript{(25)} kõrval, ja tema sõbratar on \VUgras{istmel}
\textsuperscript{(27)}, millest me näeme ainult \VUgras{seljatuge}
\textsuperscript{(26)}.

%%K t06-c2-09-1 p
9. --- Suur \VUgras{raam} \textsuperscript{(30)}, mis ripub seinal, sisaldab ühe
sugulase \VUgras{portreed} \textsuperscript{(29)}. Lauale pandud \VUgras{albumisse}
\textsuperscript{(32)} on kogutud ülejäänud pereliikmete \VUgras{fotod}
\textsuperscript{(33)}. Lapsed on seda äsja lehitsenud, sest \VUgras{toolid}
\textsuperscript{(31)} on veel laua ümber koos.

%%K t06-c2-10-1 p
10. --- Portselanist \VUgras{hiina vaas} \textsuperscript{(34)}, mis on
\VUgras{paberinoa} \textsuperscript{(35)} kõrval, kingiti Mariale tema nimepäevaks, ja
\VUgras{sirmi} \textsuperscript{(36)}, mis kaunistab toa nurka, tõi Hiinast tema onu.

%%K t06-c2-11-1 p
11. --- Rõõmustatud imeilusatest \VUgras{maalidest} \textsuperscript{(37)}, mis on
\VUgras{seina} \textsuperscript{(42)} külge riputatud, valgustatud \VUgras{laelühtrist}
\textsuperscript{(38)}, mille \VUgras{elektripirnid} \textsuperscript{(39)} õhtuti
rõõmsalt säravad, näib see salong väga meeldiv. Pehme \VUgras{vaip}
\textsuperscript{(40)}, mis on parketile laotatud, ja nii mugav \VUgras{elektrikell}
\textsuperscript{(41)} teevad selle mugavuse täielikuks.

%%K t06-tit-5 sub
\VUsaut{3.6mm}
\VUpk{10.2pt}{40}{Söögituba.}
\VUsaut{3.0mm}

%%K t06-c3-01-1 p
1. --- Õhtusöögikell on helisenud. Kogu perekond, välja arvatud Maria, on laua ümber
kogunenud. Söögituba on meeldivalt soojendatud suurepärase fajansist \VUgras{ahjuga}
\textsuperscript{(1)}, millel on peenike \VUgras{toru} \textsuperscript{(2)}.

%%K t06-c3-02-1 p
2. --- Selles toas ei ole palju mööblit. Piki seina ripub \VUgras{pingi}
\textsuperscript{(4)} kohal ilustuseks \VUgras{purskkaevuke} \textsuperscript{(3)}. Väga
lähedal on \VUgras{puhvetkapp} \textsuperscript{(5)}, millele pannakse külmad road,
magustoidutaldrikud ja mitmesugused lauatarbed, mida kohe ei vajata.

%%K t06-c3-03-1 p
3. --- Nii näeme me ülemisel riiulil \VUgras{praetud kanapoega}
\textsuperscript{(7)}, \VUgras{kala} \textsuperscript{(8)} ja \VUgras{kaussi}
\textsuperscript{(10)} \VUgras{puuviljadega} \textsuperscript{(9)}. All on siin
\VUgras{juust} \textsuperscript{(11)}, \VUgras{leib} \textsuperscript{(12)},
\VUgras{maitseainealus} \textsuperscript{(13)}, mis sisaldab kõike, mida on vaja salati
maitsestamiseks: õli, äädikat, \VUgras{soola} \textsuperscript{(14)} ja \VUgras{pipart}
\textsuperscript{(15)}. Lõpuks on alumisel riiulil \VUgras{kook} \textsuperscript{(17)}
\VUgras{sinepitoosi} \textsuperscript{(16)} kõrval.

%%K t06-c3-04-1 p
4. --- Söögitoa kell, mida nimetatakse \VUgras{seinakellaks} \textsuperscript{(20)}, on
väga erilise kujuga. See on paigutatud puitraami, mis sisaldab lisaks
\VUgras{baromeetrit} \textsuperscript{(19)} ja \VUgras{termomeetrit}
\textsuperscript{(18)}.

%%K t06-tit-6 sub
\VUsaut{4.6mm}
\VUpk{10.2pt}{40}{Õhtusöögi serveerimine.}
\VUsaut{3.0mm}

%%K t06-c3-05-1 p
5. --- Kõigile toa seintele on paigaldatud vahatatud tammest \VUgras{puitpaneelid}
\textsuperscript{(21)}. Riidest \VUgras{ukseeesriie} \textsuperscript{(22)} sulgeb
küllalt hästi ukse, mis viib verandale. Sinna läheb perekond pärast õhtusööki kohvi
jooma. Juba on \VUgras{tassid} \textsuperscript{(24)} ja \VUgras{alustassid}
\textsuperscript{(25)} valmis \VUgras{nõudekapil} \textsuperscript{(23)}.

%%K t06-c3-06-1 p
6. --- Laua kohale on lae külge kinnitatud \VUgras{rippvalgusti}
\textsuperscript{(26)}, mille väga hele lamp valgustab õhtuti kogu söögituba.

%%K t06-c3-07-1 p
7. --- Nüüd uurigem \VUgras{söögilauda} \textsuperscript{(27)} ennast. Kui perekond
sisse astus, olid lauanõud valmis. \VUgras{Laudlinale} \textsuperscript{(28)} oli iga
lauasolija kohale pandud \VUgras{taldrik} \textsuperscript{(33)}, \VUgras{salvrätik}
\textsuperscript{(29)}, \VUgras{lusikas} \textsuperscript{(34)}, \VUgras{kahvel}
\textsuperscript{(35)}, \VUgras{nuga} \textsuperscript{(36)} ja \VUgras{klaas}
\textsuperscript{(32)}. Olid ka \VUgras{pudelid} \textsuperscript{(30)} veini jaoks ja
\VUgras{karahvin} \textsuperscript{(31)} vee jaoks.

%%K t06-c3-08-1 p
8. --- \VUgras{Teener} \textsuperscript{(39)} tõi kõigepealt \VUgras{supiterriini}
\textsuperscript{(37)}, milles veel auras natuke suppi. Nüüd serveerib ta esimest
\VUgras{rooga} \textsuperscript{(38)}, liharooga. Siis pakub ta neid, mida me juba
nimetasime, ja lõpuks, pärast \VUgras{magustoiduroogi}, kasutab ta ära, et tema
peremehed on verandale läinud, et lauanõud ära koristada ja söögituba uuesti korda
seada.

%%K t06-c3-08-2 p
\textit{Märkus.} --- \textit{Toitumist puudutavad üldkasutatavad sõnad on siin antud
väga kokkuvõtlikult. Nimekirja täiendatakse kahel tabelil «Hotell» (nr.~13) ja
«Turg» (nr.~15).}

%%K t06-tit-7 sub
\VUsaut{4.6mm}
\VUpk{10.2pt}{40}{Köök.}
\VUsaut{3.0mm}

%%K t06-c4-01-1 p
1. --- Köögis, mida me hakkame vaatama läbi poolavatud ukse, näeme me maalilist
korratust. \VUgras{Kolde} \textsuperscript{(1)} ees, milles praksub \VUgras{tuli}
\textsuperscript{(3)}, on seatud \VUgras{vardapööraja} \textsuperscript{(5)}. Ilus
\VUgras{praad} \textsuperscript{(7)} on \VUgras{vardasse aetud} \textsuperscript{(6)} ja
küpseb.

%%K t06-c4-02-1 p
2. --- \VUgras{Lõõts} \textsuperscript{(8)} ja \VUgras{söelabidas}
\textsuperscript{(9)}, mida äsja kasutati, on jäänud põrandale nagu ka
\VUgras{halud} \textsuperscript{(10)}.

%%K t06-c4-03-1 p
3. --- Kamina pealne on korras; \VUgras{latern} \textsuperscript{(11)},
\VUgras{tikutoos} \textsuperscript{(13)}, milles on \VUgras{tikud}
\textsuperscript{(12)}, \VUgras{küünlajalg} \textsuperscript{(14)} ja
\VUgras{küünlahoidja} \textsuperscript{(15)} \VUgras{küünlaga} \textsuperscript{(16)} on
oma kohal. Kuid suur \VUgras{söeämber} \textsuperscript{(18)}, täis \VUgras{kivisütt}
\textsuperscript{(17)}, mille \VUgras{kokk} \textsuperscript{(23)} äsja keldris täitis,
on jäänud keset kööki.

%%K t06-c4-04-1 p
4. --- Tõesti, vaene naine peab kiirustama, et lõunasöök oleks õigeks ajaks valmis.
Nüüd on ta \VUgras{pliidi} \textsuperscript{(19)} juures ja segab \VUgras{kastet}
\textsuperscript{(24)}, mida ei tohi liialt kõrbeda lasta.

%%K t06-c4-05-1 p
5. --- \VUgras{Ahjus} \textsuperscript{(20)} küpseb kook, mille ta valmistas laste
oodeks. Pliidi kohal ripuvad \VUgras{kulp} \textsuperscript{(21)} ja
\VUgras{vahukulp} \textsuperscript{(22)}.

%%K t06-tit-8 sub
\VUsaut{4.0mm}
\VUpk{10.2pt}{40}{Nõud.}
\VUsaut{3.0mm}

%%K t06-c4-06-1 p
6. --- \VUgras{Keedunõud} \textsuperscript{(25)}, mis on toa tagaosas üles riputatud, on
väga puhtad. Ilusad vasest \VUgras{kastrulid} \textsuperscript{(26)},
\VUgras{piimapott} \textsuperscript{(27)}, \VUgras{sõel} \textsuperscript{(28)} ja
\VUgras{riiv} \textsuperscript{(29)} on hoolikalt puhastatud.

%%K t06-c4-07-1 p
7. --- \VUgras{Soolatoos} \textsuperscript{(30)} on pandud nõudekapile
\VUgras{teekannu} \textsuperscript{(31)} ja \VUgras{õlipudeli} \textsuperscript{(32)}
kõrvale. \VUgras{Sahtel} \textsuperscript{(33)} sisaldab arvatavasti söögiriistu ja
köögikahvleid.

%%K t06-c4-08-1 p
8. --- \VUgras{Luud} \textsuperscript{(34)} ja \VUgras{sulgtolmulapp}
\textsuperscript{(35)} on nurgas. Kokk kasutas äsja \VUgras{puuplokki}
\textsuperscript{(36)}, ta jättis selle akna juurde.

%%K t06-c4-09-1 p
9. --- \VUgras{Käterätik} \textsuperscript{(37)} ripub seinal, \VUgras{lehtri}
\textsuperscript{(38)} kõrval. Köögi selles osas on ära pandud \VUgras{rest}
\textsuperscript{(39)}, \VUgras{hakknuga} \textsuperscript{(40)} ja
\VUgras{lõikelaud} \textsuperscript{(41)}. Siis, hakknoa kohal, \VUgras{riiulil}
\textsuperscript{(42)}, on \VUgras{potid} \textsuperscript{(43)}, \VUgras{keedupott}
\textsuperscript{(44)} oma \VUgras{kaanega} \textsuperscript{(45)} ja \VUgras{katel}
\textsuperscript{(46)}. \VUgras{Pann} \textsuperscript{(47)} ripub lähedal.

%%K t06-c4-10-1 p
10. --- Ainult vasest \VUgras{kraan} \textsuperscript{(48)} heidab selles nurgas
sära. \VUgras{Valamu} \textsuperscript{(49)}, millele on pandud paks \VUgras{suur kann}
\textsuperscript{(50)}, on arvatavasti väga mugav oma väikese kapiga, milles hoitakse
\VUgras{äädikavaati} \textsuperscript{(51)} koos oma \VUgras{kraaniga}
\textsuperscript{(52)}, \VUgras{jäätmekasti} \textsuperscript{(53)} ja \VUgras{musti
nõusid} \textsuperscript{(54)}.

%%K t06-tit-9 sub
\VUsaut{4.0mm}
\VUpk{10.2pt}{40}{Muud kööki pandud esemed.}
\VUsaut{3.0mm}

%%K t06-c4-11-1 p
11. --- \VUgras{Tõrs} \textsuperscript{(55)}, milles pestakse \VUgras{juurvilju}
\textsuperscript{(58)}, ja malmist \VUgras{katel} \textsuperscript{(56)}, mis on pandud
\VUgras{plaatpõrandale} \textsuperscript{(57)}, on arvatavasti samuti oma koha selles
kapis.

%%K t06-c4-12-1 p
12. --- Kokal ei ole kindlasti puudust ahjudest, sest vaat et siin on lisaks, laua
nurgal, \VUgras{gaasipliit} \textsuperscript{(59)}, millel soojeneb vesi väikeses
kastrulis. \VUgras{Aur} \textsuperscript{(60)}, mis sellest väljub, näitab meile, et see
arvatavasti keeb. Tõenäoliselt kasutatakse seda vett kohvi jaoks, sest vaat et laua
teises otsas on \VUgras{kohvikann} \textsuperscript{(64)} ja \VUgras{kohviveski}
\textsuperscript{(63)}.

%%K t06-c4-13-1 p
13. --- Pliit on ühendatud \VUgras{gaasipõletiga} \textsuperscript{(62)} pika
\VUgras{kummivoolikuga} \textsuperscript{(61)}.

%%K t06-c4-14-1 p
14. --- \VUgras{Päevatööline} \textsuperscript{(66)}, kes tuleb kokka aitama, valmistab
ette juurvilju õhtusöögiks. Kui need on puhastatud ja pestud, paneb ta need
\VUgras{vaagnasse} \textsuperscript{(65)}, oodates tundi, mil neid keeta.

%%K t06-tit-10 sub
\VUsaut{4.0mm}
{\centering\fontsize{10.2pt}{10.2pt}\selectfont\textls[40]{\scshape Vannituba.} \textit{(Suplustuba.)}\par}
\VUsaut{3.0mm}

%%K t06-c5-01-1 p
1. --- Meil on jäänud teha ainult üks tagasihoidmatus, et tundma õppida korteri
peamised toad: näha vannitoa sisseseadet. See ei nõua pikka aega, sest see ei ole
suur, ja me saame kiiresti teada, et \VUgras{vannil} \textsuperscript{(1)} on hea
\VUgras{veesoojendaja} \textsuperscript{(2)}, et \VUgras{seebialus}
\textsuperscript{(3)} on suplejale käeulatuses ja et \VUgras{rätikuhoidja}
\textsuperscript{(5)} laseb kindlasti kergesti kuivada \VUgras{käterätikutel}
\textsuperscript{(4)}.

%%K t06-c5-02-1 p
2. --- Selle toa seinad on kaetud \VUgras{kahhelplaatidega} \textsuperscript{(6)}, mis
võimaldasid paigaldada \VUgras{duši} \textsuperscript{(7)}, kartmata, et vesi, mis selle
kasutamisel tagasi pritsib, müüre kahjustab. Tsingist \VUgras{jalavann}
\textsuperscript{(8)}, mis on jalasuplusteks, teeb sisustuse täielikuks.

\VUsaut{6.0mm}
\VUfilet{22mm}
